% !TeX program = xelatex
%% =====================================================================
%%  [เฉลย] ใบกิจกรรมในชั้นเรียน (Worksheet Solution) บทที่ 3 : ความสัมพันธ์และฟังก์ชัน --- สัปดาห์ที่ 7
%%  เอกสารเฉลยสำหรับผู้สอน: 2 หน้าเต็ม (ฉบับสมบูรณ์)
%% =====================================================================
\documentclass[a4paper,12pt]{article}
\usepackage{fontspec}
\usepackage{amsmath,amssymb}
\usepackage[margin=1.2cm,top=1.0cm,bottom=1.0cm]{geometry}
\usepackage{enumitem}
\usepackage{array}
\usepackage{xcolor}
\usepackage{tikz}

\XeTeXlinebreaklocale "th_TH"
\XeTeXlinebreakskip = 0pt plus 1pt
\setmainfont[Script=Thai,Scale=1.15]{TH Sarabun New}
\definecolor{navy}{HTML}{1F3A93}
\definecolor{mint}{HTML}{E8F5E9}
\definecolor{lightnavy}{HTML}{EAEEF7}
\definecolor{ans}{HTML}{C0392B} % สีเฉลยเน้นชัดเจน
\newcommand{\ans}[1]{\textbf{\color{ans}#1}}

\renewcommand{\arraystretch}{1.15}
\setlist{topsep=1pt,itemsep=2pt,parsep=0pt}
\pagestyle{empty}

\begin{document}

%% ---------- หัวกระดาษหน้า 1 ----------
\noindent
\begin{minipage}[c]{0.78\textwidth}
  {\Large\bfseries\color{navy} [เฉลย] ใบกิจกรรมในชั้นเรียน บทที่ 3 : ความสัมพันธ์และฟังก์ชัน (สัปดาห์ที่ 7)}\\[0.1em]
  {\normalsize รายวิชา 4091606 คณิตศาสตร์สำหรับคอมพิวเตอร์}\\[0.5em]
  {\bfseries\color{ans} ฉบับเฉลยและแนวทางคำตอบสำหรับผู้สอน (Master Answer Key)}
\end{minipage}\hfill
\begin{minipage}[c]{0.2\textwidth}
  \centering
  \fbox{\parbox{2.8cm}{\centering\bfseries\color{navy}เฉลยประจำสัปดาห์\\ Week 07 KEY}}
\end{minipage}
\par\vspace{0.3em}\hrule\vspace{0.4em}

\noindent\textbf{คำชี้แจง:} จงแสดงวิธีคิดและเติมคำตอบลงในช่องว่างให้สมบูรณ์ เพื่อร่วมอภิปรายและประเมินผลในชั้นเรียน
\vspace{0.3em}

%% ---------- กิจกรรม 1 ----------
\noindent\textbf{กิจกรรมที่ 1 : พื้นฐานคู่อันดับ ผลคูณคาร์ทีเซียน และการหาจำนวนความสัมพันธ์} \hfill \textit{(6 นาที)}
\begin{enumerate}[label=\arabic*)]
  \item กำหนด $A = \{2, 5, 8\}$ และ $B = \{x, y\}$ จงหา:
        \begin{enumerate}[label=(\alph*),itemsep=2pt]
          \item จำนวนสมาชิก $|A \times B| =$ \ans{6 \quad ($|A| \cdot |B| = 3 \times 2 = 6$)}
          \item จำนวนความสัมพันธ์ทั้งหมดจาก $A$ ไป $B = 2^{|A \times B|} =$ \ans{$2^6 = 64$ \text{ความสัมพันธ์}}
          \item จำนวนความสัมพันธ์ทวิภาคทั้งหมดบนเซต $A = 2^{|A|^2} =$ \ans{$2^{3^2} = 2^9 = 512$ \text{ความสัมพันธ์}}
        \end{enumerate}
  \item กำหนดเซตของจำนวนเฉพาะ $S = \{2, 3, 5, 7\}$ และความสัมพันธ์ $R = \{(a, b) \in S \times S \mid a < b\}$\\
        จงแจกแจงสมาชิกของ $R = $ \ans{\{(2, 3), (2, 5), (2, 7), (3, 5), (3, 7), (5, 7)\}}
  \item จงหาโดเมนและเรนจ์ของ $R$ ในข้อ 2):
        \quad $\operatorname{Dom}(R) = $ \ans{\{2, 3, 5\}}
        \quad $\operatorname{Ran}(R) = $ \ans{\{3, 5, 7\}}
\end{enumerate}
\vspace{0.6em}

%% ---------- กิจกรรม 2 ----------
\noindent\textbf{กิจกรรมที่ 2 : ความสัมพันธ์ผกผันและการดำเนินการระหว่างความสัมพันธ์} \hfill \textit{(6 นาที)}
\begin{enumerate}[label=\arabic*)]
  \item กำหนด $R_1 = \{(a, 10), (b, 20), (c, 10)\}$ จากเซต $\{a, b, c\}$ ไปยังเซต $\{10, 20\}$\\
        จงหาความสัมพันธ์ผกผัน $R_1^{-1} = $ \ans{\{(10, a), (20, b), (10, c)\}}
  \item กำหนด $R_2 = \{(a, 20), (b, 20), (c, 10)\}$ จงหา:
        \begin{enumerate}[label=(\alph*),itemsep=2pt]
          \item $R_1 \cup R_2 =$ \ans{\{(a, 10), (a, 20), (b, 20), (c, 10)\}}
          \item $R_1 \cap R_2 =$ \ans{\{(b, 20), (c, 10)\}}
          \item $R_1 - R_2 =$ \ans{\{(a, 10)\}}
        \end{enumerate}
\end{enumerate}
\vspace{0.6em}

%% ---------- กิจกรรม 3 ----------
\noindent\textbf{กิจกรรมที่ 3 : การตรวจสอบสมบัติของความสัมพันธ์ทวิภาค (Reflexive, Symmetric, Transitive)} \hfill \textit{(8 นาที)}\\
กำหนดเซต $A = \{p, q, r\}$ จงตรวจสอบสมบัติของความสัมพันธ์ต่อไปนี้ (ใส่เครื่องหมาย $\checkmark$ หรือ $\times$):
\begin{center}
\renewcommand{\arraystretch}{1.3}
\small
\begin{tabular}{|l|c|c|c|c|}
\hline
\textbf{ความสัมพันธ์บน $A$} & \textbf{สะท้อน (Refl.)} & \textbf{สมมาตร (Sym.)} & \textbf{ปฏิสมมาตร (Antisym.)} & \textbf{ถ่ายทอด (Trans.)} \\ \hline
$R_A = \{(p,p), (q,q), (r,r), (p,q), (q,p)\}$ & \ans{$\checkmark$} & \ans{$\checkmark$} & \ans{$\times$} & \ans{$\checkmark$} \\ \hline
$R_B = \{(p,q), (q,r), (p,r)\}$ & \ans{$\times$} & \ans{$\times$} & \ans{$\checkmark$} & \ans{$\checkmark$} \\ \hline
$R_C = \{(p,p), (q,q), (r,r), (p,q), (q,r), (p,r)\}$ & \ans{$\checkmark$} & \ans{$\times$} & \ans{$\checkmark$} & \ans{$\checkmark$} \\ \hline
$R_D = \{(p,q), (q,p)\}$ & \ans{$\times$} & \ans{$\checkmark$} & \ans{$\times$} & \ans{$\times$} \\ \hline
\end{tabular}
\end{center}

\clearpage
%% ====================== หน้าที่ 2 ======================
\noindent
\begin{minipage}[c]{\textwidth}
  {\large\bfseries\color{navy} [เฉลย] ใบกิจกรรมในชั้นเรียน บทที่ 3 : ความสัมพันธ์และฟังก์ชัน (สัปดาห์ที่ 7) --- หน้า 2}\\[0.25em]
  {\bfseries\color{ans} เฉลยกิจกรรมที่ 4 และ 5 (Master Answer Key)}
\end{minipage}
\par\vspace{0.3em}\hrule\vspace{0.4em}

\noindent\textbf{คำชี้แจง:} จงแสดงวิธีคิดและหาคำตอบตามเนื้อหาบทเรียน (ทำต่อเนื่องจากหน้า 1)
\vspace{0.5em}

%% ---------- กิจกรรม 4 ----------
\noindent\textbf{กิจกรรมที่ 4 : การประกอบความสัมพันธ์และเมทริกซ์ 0-1 (Composite Relations \& 0-1 Matrix)} \hfill \textit{(10 นาที)}\\
กำหนด $A = \{x, y, z\}, B = \{1, 2, 3\}, C = \{a, b\}$ โดย $R = \{(x, 1), (x, 2), (y, 3), (z, 1)\}$ และ $S = \{(1, a), (2, b), (3, a)\}$
\begin{enumerate}[label=\arabic*)]
  \item จงหาความสัมพันธ์ประกอบ $S \circ R$ (จาก $A$ ไป $C$ โดยผ่าน $B$):\\[0.3em]
        $S \circ R =$ \ans{\{(x, a), (x, b), (y, a), (z, a)\}}
  \item เขียนเมทริกซ์บูลีน $M_R$ (ขนาด $3 \times 3$ ของแถว $x,y,z$ และคอลัมน์ $1,2,3$) และ $M_S$ (แถว $1,2,3$ และคอลัมน์ $a,b$):
        \[
        M_R = \begin{bmatrix}
        \ans{1} & \ans{1} & \ans{0} \\[0.3em]
        \ans{0} & \ans{0} & \ans{1} \\[0.3em]
        \ans{1} & \ans{0} & \ans{0}
        \end{bmatrix}, \quad
        M_S = \begin{bmatrix}
        \ans{1} & \ans{0} \\[0.3em]
        \ans{0} & \ans{1} \\[0.3em]
        \ans{1} & \ans{0}
        \end{bmatrix}
        \quad \implies M_{S \circ R} = M_R \odot M_S = \begin{bmatrix}
        \ans{1} & \ans{1} \\[0.3em]
        \ans{1} & \ans{0} \\[0.3em]
        \ans{1} & \ans{0}
        \end{bmatrix}
        \]
\end{enumerate}
\vspace{0.8em}

%% ---------- กิจกรรม 5 ----------
\noindent\textbf{กิจกรรมที่ 5 : การปิดสมบัติของความสัมพันธ์ (Closures of Relations)} \hfill \textit{(10 นาที)}
\begin{enumerate}[label=\arabic*)]
  \item กำหนดเซต $A = \{a, b, c, d\}$ และความสัมพันธ์ $R_1 = \{(a, b), (b, c), (c, d)\}$ บนเซต $A$ จงหา:
        \begin{enumerate}[label=(\alph*),itemsep=3pt]
          \item การปิดสมบัติสะท้อน (Reflexive Closure): $r(R_1) =$ \ans{\{(a, b), (b, c), (c, d), (a, a), (b, b), (c, c), (d, d)\}}
          \item การปิดสมบัติสมมาตร (Symmetric Closure): $s(R_1) =$ \ans{\{(a, b), (b, c), (c, d), (b, a), (c, b), (d, c)\}}
          \item การปิดสมบัติถ่ายทอด (Transitive Closure): $t(R_1) =$ \ans{\{(a, b), (b, c), (c, d), (a, c), (b, d), (a, d)\}}
        \end{enumerate}
  \item กำหนดเซต $B = \{2, 4, 6\}$ และความสัมพันธ์ $R_2 = \{(2, 2), (2, 4), (4, 6)\}$ บนเซต $B$ จงหา:
        \begin{enumerate}[label=(\alph*),itemsep=3pt]
          \item การปิดสมบัติสะท้อน (Reflexive Closure): $r(R_2) =$ \ans{\{(2, 2), (2, 4), (4, 6), (4, 4), (6, 6)\}}
          \item การปิดสมบัติสมมาตร (Symmetric Closure): $s(R_2) =$ \ans{\{(2, 2), (2, 4), (4, 6), (4, 2), (6, 4)\}}
          \item การปิดสมบัติถ่ายทอด (Transitive Closure): $t(R_2) =$ \ans{\{(2, 2), (2, 4), (4, 6), (2, 6)\}}
        \end{enumerate}
\end{enumerate}

\end{document}
